
\subsection{散射数据与矩阵问题}
本节先规定逆问题, 再说明正向散射产生的矩阵满足这些条件.
取满足共轭、换叶约化的演化反射数据 $\rho(z,t),\widetilde\rho(z,t)$,
以及有限简单离散数据 $\zeta_j,c_j(0)$, $j=1,\ldots,m$, $m=2N$.
各离散点在 $\Sigma$ 外, 透射系数在普通连续谱上非零,
并取 \eqref{eq:C-timephase} 的 $C_j(x,t)$.
共振处要求 \eqref{eq:resonance-reflection} 的反射条件.
谱数据的相位约束在本节末尾另行说明.

将 $\Sigma$ 定向为 $D_+$ 位于左侧: 实轴外部两段向右、内部区间向左,
上半圆顺时针、下半圆逆时针. 下标 $+$、$-$ 表示左、右非切向边界值.
普通谱弧删去 $0,\pm q_0,\pm\ii q_0,z_a^\pm$.

\begin{definition}[非零背景矩阵 RHP]\label{def:RHP}\label{thm:RHP}
求关于 $z$ 分片亚纯的二阶矩阵 $M(x,t,z)$, 满足下列条件.
在 $D_+$ 中只有第二列可在 $\zeta_j$ 有简单极点,
在 $D_-$ 中只有第一列可在 $\conj{\zeta_j}$ 有简单极点, 留数由 \eqref{eq:resplus}--\eqref{eq:resminus} 给定.
要求普通谱弧上具有局部连续的非切向边界值, 跳跃为
\begin{equation}\label{eq:jump}
M_+=M_-J,\qquad
J=\begin{pmatrix}
1&\widetilde\rho(z,t)e^{-2\ii\lambda x}\\
-\rho(z,t)e^{2\ii\lambda x}&1-\rho(z,t)\widetilde\rho(z,t)
\end{pmatrix},\qquad \det J=1.
\end{equation}
两处归一化为
\begin{align}
M(z)&=I+\dfrac{M^{(1)}}{z}+O(z^{-2}),\qquad z\to\infty,\label{eq:Minfty}\\
M(z)&=-\dfrac{\ii q_0}{z}\sigma_1+O(1),\qquad z\to0.\label{eq:Mzero}
\end{align}
要求约化关系
\begin{equation}\label{eq:Msym}
M(z)=\sigma_2\conj{M(\conj z)}\sigma_2,\qquad
M(h(z))=M(z)\dfrac{\ii z}{q_0}\sigma_1.
\end{equation}
除上述指定主部外, 矩阵在各有限特殊点局部有界,
其扇区边界值与相邻谱弧上的跳跃相容.
在两个共振原像处要求两侧具有共同连续值.
\end{definition}
这里的局部要求是 RHP 解类的一部分.
Jost 列在分支点连续不自动保证拼接矩阵满足这一要求;
对独立数据, 本定义也不预先断言解存在.

\begin{proposition}[由正向散射构造 RHP]\label{prop:jump}
设三场满足第 \ref{ch:scattering} 节的正向演化条件,
其透射系数在普通连续谱上非零且离散零点有限、简单.
以 Jost 解析列构造
\begin{equation}\label{eq:Mdef}
M(z)=\begin{cases}
M_+(z)=\left(\mu_{-,1},\dfrac{\mu_{+,2}}{s_{22}}\right),&z\in D_+,\\[6pt]
M_-(z)=\left(\dfrac{\mu_{+,1}}{s_{11}},\mu_{-,2}\right),&z\in D_-.
\end{cases}
\end{equation}
若该拼接矩阵具有定义 \ref{def:RHP} 规定的有限特殊点边界行为,
则它是相应 RHP 的解.
\end{proposition}
\begin{proof}
解析性与列极点由第 \ref{ch:direct} 节的结果得到. 去除散射关系中的指数, 有
\[
\mu_{+,1}=s_{11}\mu_{-,1}+s_{21}e^{2\ii\lambda x}\mu_{-,2},\qquad
\mu_{+,2}=s_{12}e^{-2\ii\lambda x}\mu_{-,1}+s_{22}\mu_{-,2}.
\]
第一式给出 $M_{+,1}=M_{-,1}-\rho e^{2\ii\lambda x}M_{-,2}$, 代入第二式得到
\[
M_{+,2}=\widetilde\rho e^{-2\ii\lambda x}M_{-,1}+(1-\rho\widetilde\rho)M_{-,2},
\]
即 \eqref{eq:jump}. Wronskian 表达式给出
\[
\det M_+=\dfrac{\det(\mu_{-,1},\mu_{+,2})}{s_{22}}=\gamma,
\]
另一域同理. 第 \ref{ch:asymptotics} 节的展开确定两处归一化. 共轭约化交换两列和两域; 换叶约化将 $\mu_{-,1}$ 与 $\mu_{-,2}$、$\dfrac{\mu_{+,2}}{s_{22}}$ 与 $\dfrac{\mu_{+,1}}{s_{11}}$ 配对, 因而得到 \eqref{eq:Msym}.
\end{proof}


\begin{lemma}[RHP 行列式]\label{lem:RHP-determinant}
任一属于定义 \ref{def:RHP} 解类的矩阵满足
\begin{equation}\label{eq:detM}
\det M(z)=\gamma(z)=1+\dfrac{q_0^2}{z^2}.
\end{equation}
\end{lemma}
\begin{proof}
由 $\det J=1$, 行列式在普通谱弧上没有跳跃.
在离散极点, 奇异列的留数是正则列的标量倍数, 所以行列式的留数为零.
有限特殊点的局部有界性使它们对行列式可去.
两处归一化因而给出 $\det M=1+d/z+q_0^2/z^2$.
另一方面, 换叶约化给出
\[
\det M(h(z))=\dfrac{z^2}{q_0^2}\det M(z).
\]
比较两边的一次项, 得 $-d=d$, 故 $d=0$.
\end{proof}
\subsection{含反射的 Cauchy 表示}
记 $e_1,e_2$ 为标准基, 置
\begin{equation}\label{eq:Gdefinition}
G=I-J^{-1}=\begin{pmatrix}
\rho\widetilde\rho&\widetilde\rho e^{-2\ii\lambda x}\\
-\rho e^{2\ii\lambda x}&0
\end{pmatrix},\qquad M_+-M_-=M_+G.
\end{equation}
以下表示按照先消去奇性、再应用 Plemelj 公式的顺序建立\cite{ZhangFan2020,Fan2022}. 它适用于满足所列条件的 RHP 解.

\begin{proposition}[含反射的留数积分表示]\label{prop:cauchy-representation}
设定义 \ref{def:RHP} 的 RHP 有解. 假设减去离散主部及 $z=0$ 主部后, 矩阵在各有限特殊点有有界、相容的扇区边界值; 跳跃密度 $W=M_+G$ 在普通弧上局部 Hölder 连续, 且 $W,sW\in L^1(\Sigma)$. 删去特殊点的小圆弧时, 假设其上的正则化矩阵一致有界. 则
\begin{equation}\label{eq:general-cauchy}
\begin{aligned}
M(z)={}&I-\dfrac{\ii q_0}{z}\sigma_1
+\sum\limits_{j=1}^{m}\left(
\dfrac{C_j\mathbf v_je_2^T}{z-\zeta_j}
-\dfrac{\conj{C_j}\mathbf u_je_1^T}{z-\conj{\zeta_j}}\right)\\
&+\dfrac{1}{2\pi\ii}\int\limits_\Sigma\dfrac{M_+(s)G(s)}{s-z}\,\dd s,
\qquad z\in\mathbb C\mathbin{\backslash}\Sigma,
\end{aligned}
\end{equation}
其中 $\mathbf v_j=M_1(\zeta_j)$、$\mathbf u_j=M_2(\conj{\zeta_j})$ 是各点的正则列. 它们满足积分--代数关系
\begin{align}
\mathbf v_j={}&e_1-\dfrac{\ii q_0}{\zeta_j}e_2
-\sum\limits_{k=1}^{m}\dfrac{\conj{C_k}\mathbf u_k}{\zeta_j-\conj{\zeta_k}}
+\dfrac{1}{2\pi\ii}\int\limits_\Sigma\dfrac{W_1(s)}{s-\zeta_j}\,\dd s,
\label{eq:general-finite1}\\
\mathbf u_j={}&e_2-\dfrac{\ii q_0}{\conj{\zeta_j}}e_1
+\sum\limits_{k=1}^{m}\dfrac{C_k\mathbf v_k}{\conj{\zeta_j}-\zeta_k}
+\dfrac{1}{2\pi\ii}\int\limits_\Sigma\dfrac{W_2(s)}{s-\conj{\zeta_j}}\,\dd s.
\label{eq:general-finite2}
\end{align}
按 $q=\ii M^{(1)}_{12}$ 定义重构光场, 则
\begin{equation}\label{eq:qgeneral}
q=q_0+\ii\sum\limits_{j=1}^{m}C_j(\mathbf v_j)_1
-\dfrac{1}{2\pi}\int\limits_\Sigma W_{12}(s)\,\dd s.
\end{equation}
\end{proposition}
\begin{proof}
第 \ref{ch:discrete} 节的留数条件给出
\[
\Res_{z=\zeta_j}M=C_j\mathbf v_je_2^T,\qquad
\Res_{z=\conj{\zeta_j}}M=-\conj{C_j}\mathbf u_je_1^T.
\]
从 $M$ 减去 $I-\dfrac{\ii q_0}{z}\sigma_1$ 及上述各主部, 所得分片解析矩阵 $F_0$ 在无穷远为 $O(z^{-1})$, 而 $F_{0,+}-F_{0,-}=W$. 对删去特殊点小圆盘并截断于大圆的两域分别使用 Cauchy 公式. 左侧为正向边界, 将两域的公共边界合并得到 $W$; 局部一致有界性保证小圆贡献消失, 大圆上的 $O(z^{-1})$ 使其贡献消失. 因而
\[
F_0(z)=\dfrac{1}{2\pi\ii}\int\limits_\Sigma\dfrac{W(s)}{s-z}\,\dd s,
\]
得到 \eqref{eq:general-cauchy}. 正号由 \eqref{eq:Gdefinition} 与轮廓定向共同确定.

在第一列取 $z=\zeta_j$, 在第二列取 $z=\conj{\zeta_j}$, 各自本列在该点没有离散极点, 得 \eqref{eq:general-finite1}--\eqref{eq:general-finite2}. 沿避开轮廓的闭角域令 $z\to\infty$, 使用
\[
\dfrac{1}{s-z}=-\dfrac{1}{z}+\dfrac{s}{z(s-z)},\qquad
\left|\dfrac{s}{z(s-z)}\right|\le\dfrac{C|s|}{|z|^2},
\]
及 $sW\in L^1$ 控制余项, 得
\begin{equation}\label{eq:M1-general}
M^{(1)}=-\ii q_0\sigma_1
+\sum\limits_{j=1}^{m}(C_j\mathbf v_je_2^T-\conj{C_j}\mathbf u_je_1^T)
-\dfrac{1}{2\pi\ii}\int\limits_\Sigma W(s)\,\dd s.
\end{equation}
读取 $(1,2)$ 元的系数并乘以 $\ii$, 即得 \eqref{eq:qgeneral}. 对直接散射构造的 $M$, 此定义与原光场一致, 见下述重构命题.
\end{proof}

\subsection{光场与介质的恢复公式}
\begin{proposition}[正向散射的恢复恒等式]\label{prop:reconstruction}
直接散射构造的矩阵 $M$ 满足
\begin{equation}\label{eq:reconq}
Q(q)=\dfrac{\ii}{2}[\sg,M^{(1)}],\qquad q=\ii M^{(1)}_{12}.
\end{equation}
若满足命题 \ref{prop:resonance} 的共振相容条件, 则 $M$ 在两个共振原像的空间归一化值可逆, 且
\begin{equation}\label{eq:RfromM}
R=M(x,t,z_a)\dfrac{\eta_0\lambda(z_a)}{a}\sg M(x,t,z_a)^{-1},
\qquad z_a=z_a^\pm.
\end{equation}
两式给出同一个介质矩阵. 原物理场由
\begin{equation}\label{eq:original-fields}
\widetilde q=e^{-\ii\Omega t}q,\qquad
\widetilde p=-e^{-\ii\Omega t}R_{12},\qquad \eta=R_{11}
\end{equation}
恢复.
\end{proposition}
\begin{proof}
由 \eqref{eq:mu-highenergy} 与 \eqref{eq:transmission-highenergy}, 在 $D_+$ 中第二列除以 $s_{22}$ 后, 其对角系数从 $-\ii J_+$ 变为 $-\ii(J_++J_{\rm tot})=-\ii J_-$. 另一域相同, 故
\[
M^{(1)}=-\ii B(q)+\ii J_-\sg,
\]
得到 \eqref{eq:reconq}. 这一关系也可由空间方程
\begin{equation}\label{eq:Mspace}
M_x=-\ii k\sg M+QM+\ii\lambda M\sg
\end{equation}
比较无穷远的常数项得到: 第 \ref{ch:asymptotics} 节的积分估计还给出 $M_x=O(z^{-1})$, 而 $k,\lambda\sim\dfrac{z}{2}$.

在共振点, $S$ 对角, 两种拼接均为 $M=\mu_-$. 命题 \ref{prop:resonance} 的证明给出
\[
\chi_-^{-1}R\chi_-=\dfrac{\eta_0\lambda(z_a)}{a}\sg.
\]
空间指数与 $\sg$ 对易, 因而得到 \eqref{eq:RfromM}; 这里使用空间归一化值, 不对奇异的传播指数求值. $\gamma(z_a)>0$ 保证可逆. 换叶时 $\lambda$ 变号, 同时 $\sigma_1\sg\sigma_1=-\sg$, 两次变号相消, 故两原像的结果一致.
\end{proof}


\begin{lemma}[介质重构的约化相容性]\label{lem:medium-reduction}
设 $M$ 是定义 \ref{def:RHP} 的一个解.
用 \eqref{eq:RfromM} 定义 $R$, 则共振处逆矩阵存在, 两个原像给出同一个 $R$.
该矩阵为 Hermitian 迹零矩阵, 且
\begin{equation}\label{eq:reconstructed-bloch}
R^2=\eta_0^2\left(1+\dfrac{q_0^2}{a^2}\right)I.
\end{equation}
\end{lemma}
\begin{proof}
引理 \ref{lem:RHP-determinant} 给出 $\det M(z_a)=\gamma(z_a)>0$.
换叶约化中 $\lambda(h(z_a))=-\lambda(z_a)$,
而 $\sigma_1\sg\sigma_1=-\sg$; 代入 \eqref{eq:RfromM}, 两次变号抵消, 得到同一个矩阵.

共振处两侧有共同值, 实轴上的共轭约化使它具有形式
\[
M(z_a)=\begin{pmatrix}u&v\\-\conj v&\conj u\end{pmatrix}.
\]
因此 $M(z_a)^\dagger M(z_a)=\gamma(z_a)I$.
式 \eqref{eq:RfromM} 是实 Hermitian 对角矩阵的酉相似变换,
从而保持 Hermitian 性与迹零. 其平方为
$\dfrac{\eta_0^2\lambda(z_a)^2}{a^2}I$;
再用 $\lambda(z_a)^2=a^2+q_0^2$, 得到 \eqref{eq:reconstructed-bloch}.
\end{proof}
引理 \ref{lem:medium-reduction} 是对任一满足所列条件的矩阵解的代数结论.
要由独立数据得到 HMB 解, 还需验证参数正则性、空间端态及原方程.
第 \ref{ch:verification} 节对指定纯虚单轨道完成这些步骤.
\subsection{迹公式与相位条件}
本小节各谱数据取同一时刻. 相位条件由 $z=0$ 处的归一化确定. 透射系数同时包含离散零点和连续谱反射的贡献. 以下标量分解采用非零边界逆散射中的 Cauchy 积分方法\cite{Fan2022,ZhangFan2020}, 积分轮廓取 HMB 的实轴与圆周.

\begin{proposition}[迹公式与同相约束]\label{prop:theta}
设透射系数在普通轮廓上非零, 只有上述有限简单离散零点. 取
$\ell(s)=\log(1-\rho(s)\widetilde\rho(s))$ 的零指标分支, 设其在普通弧上局部 Hölder 连续, 在 $0,\infty$ 的极限为零, 且 $\ell,\dfrac{\ell(s)}{s}\in L^1(\Sigma)$. 假设剔除零点后两域中的透射函数具有与乘积关系相容的单值对数, 这些对数与下述 Cauchy 积分在各有限特殊点的边界值有界并局部匹配. 则
\begin{equation}\label{eq:trace}
s_{22}(z)=\prod_{j=1}^{m}\dfrac{z-\zeta_j}{z-\conj{\zeta_j}}
\exp\left\{-\dfrac{1}{2\pi\ii}\int\limits_\Sigma\dfrac{\ell(s)}{s-z}\,\dd s\right\},\qquad z\in D_+.
\end{equation}
$s_{11}$ 由共轭约化确定. 相同左右背景要求
\begin{equation}\label{eq:theta-condition}
\prod_{j=1}^{m}\dfrac{\zeta_j}{\conj{\zeta_j}}
\exp\left\{-\dfrac{1}{2\pi\ii}\int\limits_\Sigma\dfrac{\ell(s)}{s}\,\dd s\right\}=1.
\end{equation}
特别地, 无反射时 $s_{22}$ 等于 \eqref{eq:trace} 中的有理乘积, 且 $\displaystyle\prod_{j=1}^{m}\dfrac{\zeta_j}{\conj{\zeta_j}}=1$.
\end{proposition}
\begin{proof}
记 $B_0(z)=\displaystyle\prod_{j=1}^{m}\dfrac{z-\zeta_j}{z-\conj{\zeta_j}}$, 定义
$\beta_+=\dfrac{s_{22}}{B_0}$、$\beta_-=s_{11}B_0$. 它们在相应域解析且无零点, 无穷远归一化为 1. 由 $\det S=1$,
\[
\beta_+\beta_-=s_{11}s_{22}=\dfrac{1}{1-\rho\widetilde\rho}.
\]
相容的对数分支满足 $\log\beta_++\log\beta_-=-\ell$, 因而分片解析函数 $f_+=\log\beta_+$、$f_-=-\log\beta_-$ 的跳跃为 $f_+-f_-=-\ell$. Plemelj 公式及归一化给出
\[
f(z)=-\dfrac{1}{2\pi\ii}\int\limits_\Sigma\dfrac{\ell(s)}{s-z}\,\dd s.
\]
差函数在特殊点可去且在无穷远为零, 因而没有附加整函数项. 代回 $\beta_+$ 即得 \eqref{eq:trace}. 为取零点极限, 沿 $D_+$ 内的 $z=-\ii\varepsilon$, $\varepsilon\downarrow0$. 实轴上 $\left|\dfrac{s}{s+\ii\varepsilon}\right|\le1$, 圆周与零点分离, 故 $\dfrac{\ell(s)}{s-z}$ 受 $C\left|\dfrac{\ell(s)}{s}\right|$ 控制. 支配收敛结合 $s_{22}(0)=1$ 给出 \eqref{eq:theta-condition}.
\end{proof}

每个完整四点轨道取外部上半平面的代表 $\zeta_n$, 则它的两个 $D_+$ 零点对无反射相位乘积的贡献为 $\left(\dfrac{\zeta_n}{\conj{\zeta_n}}\right)^2$. 因而同相约束化为
\[
4\sum\limits_{n=1}^{N}\arg\zeta_n=0\pmod{2\pi}.
\]
纯虚代表逐轨道满足这一条件.
